\documentclass[11pt,a4paper]{article}
\usepackage[T1]{fontenc}
\usepackage[utf8]{inputenc}
\usepackage[brazil]{babel}
\usepackage{lmodern}
\usepackage{amsmath,amssymb,mathtools}
\usepackage{geometry,microtype,booktabs,tabularx}
\geometry{margin=2.2cm}
\newtheorem{theorem}{Teorema}[section]
\newtheorem{lemma}[theorem]{Lema}
\newtheorem{proposition}[theorem]{Proposição}
\newcommand{\X}{\mathcal X}
\newcommand{\Mmot}{\mathcal M}
\title{Obstruções Mínimas de Medida Positiva na Dinâmica Multilinear Projetada de 3-SAT\\\large Uma armadilha hiperarbórea de seis cláusulas e sua ocorrência com densidade linear em fórmulas aleatórias esparsas}
\author{Thiago Carvalho\\Pesquisador Independente\\\texttt{thiagocarvalho.dba@gmail.com}}
\date{Versão para análise -- derivada da fonte auditada CLG-R v4.0.5 (setembro de 2026)}
\begin{document}
\maketitle
\begin{abstract}
Estudamos a dinâmica de gradiente projetado para a extensão multilinear natural de 3-SAT no hipercubo $[-1,1]^N$. A questão determinística central é identificar o menor componente conexo, linear, 3-uniforme e Berge-acíclico capaz de sustentar uma bacia de medida de Lebesgue positiva cujos limites, após arredondamento Booleano, violem alguma cláusula. Provamos que o mínimo é seis cláusulas. O limite inferior para todos os componentes com no máximo cinco cláusulas é obtido sem argumentos de variedade estável: cláusulas positivas forçam variáveis de grau um, cada uma requer uma cláusula nula restauradora distinta, e a estrutura arbórea resultante contém um ramo unitário com uma coordenada-folha globalmente monótona. A convergência a um equilíbrio de energia positiva então força a condição inicial sobre um hiperplano de bordo, portanto um conjunto nulo. Para o limite superior damos um motivo explícito M6 com família de equilíbrios de dimensão nove em energia um e uma caixa aberta de captura com massa certificada $2^{-53}$. Em seguida contamos cópias isoladas assinadas de M6 no ensemble uniforme esparso de 3-SAT. Sua densidade esperada converge para
\[
d(\alpha)=\frac{729}{64}\alpha^6e^{-39\alpha}.
\]
Um argumento de diferenças limitadas por troca de cláusula e um cálculo binomial condicional fornecem uma cota inferior exponencialmente de alta probabilidade para a densidade residual final. Em particular,
\[
\liminf_{N\to\infty}\mathbb E\rho_{N,\mathrm{mult}}\ge \frac{729}{2^{59}}\alpha^5e^{-39\alpha}>0,
\]
e a cota conservadora com constante $729/2^{60}$ vale com probabilidade tendendo a um. O resultado trata da dinâmica desta representação contínua e não implica nada sobre o problema $P$ versus $NP$.
\end{abstract}
\noindent\textbf{Palavras-chave:} sistemas dinâmicos projetados; 3-SAT; extensão multilinear; hipergrafos aleatórios; desigualdade de Kurdyka--\L{}ojasiewicz; atratores espúrios.

\section{Introdução}
Um objetivo Booleano pode ser estendido dos vértices de um hipercubo para seu interior de muitas formas inequivalentes. Mesmo quando duas extensões coincidem exatamente em todas as atribuições Booleanas, sua geometria crítica contínua e sua dinâmica de gradiente projetado podem diferir substancialmente. Isolamos aqui um mecanismo concreto na representação multilinear de 3-SAT: um pequeno componente acíclico pode possuir uma bacia aberta convergindo para uma família não estrita de equilíbrios com energia positiva.

A contribuição tem duas partes. Primeiro, na classe HT de hipergrafos de cláusulas conexos, lineares, 3-uniformes e Berge-acíclicos, determinamos o número mínimo exato de cláusulas necessário para uma bacia ruim de medida positiva. Segundo, mostramos que o motivo mínimo ocorre com densidade linear em um ensemble aleatório esparso, de modo que a obstrução não é apenas um exemplo patológico finito.

\section{Dinâmica multilinear projetada de 3-SAT}
Seja $F$ uma fórmula 3-CNF com $M$ cláusulas em $N$ variáveis. Para cada cláusula $c$, seja $\sigma^{(c)}\in\{-1,0,+1\}^N$ seu vetor de polaridades, com exatamente três entradas não nulas. O potencial multilinear natural de violação é
\begin{equation}
\Phi_{\rm mult}(x)=\sum_{c=1}^M\prod_{j\in c}\frac{1-\sigma_j^{(c)}x_j}{2},\qquad x\in\X=[-1,1]^N.
\end{equation}
Nos vértices Booleanos, ele coincide com o número de cláusulas violadas. Consideramos o fluxo
\begin{equation}
\dot x=\Pi_{T_\X(x)}(-\nabla\Phi_{\rm mult}(x)).
\end{equation}
Para argumentos locais, orientamos cada variável e usamos coordenadas-fator $z_v\in[0,1]$, de modo que cada fator literal seja $z_v$ ou $1-z_v$. Se $z=(1\pm x)/2$, então, no tempo original,
\begin{equation}
\dot z=\frac14\Pi_{T_{[0,1]^n}(z)}(-\nabla_z\Phi).
\end{equation}
A reparametrização temporal positiva não altera órbitas, equilíbrios ou bacias.

\subsection{Convergência pontual}
Defina $G(z)=\Phi(z)+\delta_{[0,1]^n}(z)$. Como $\Phi$ é polinomial e a caixa é semialgébrica compacta, $G$ é própria, semicontínua inferior e semialgébrica. A decomposição de Moreau fornece
\[
-\nabla\Phi=d+n,\qquad d=\Pi_T(-\nabla\Phi),\qquad n\in N_X,\qquad \langle d,n\rangle=0,
\]
logo
\[
-d\in\partial G(z),\qquad \|d\|=\operatorname{dist}(0,\partial G(z)),
\]
e, ao longo de toda solução absolutamente contínua,
\begin{equation}
\frac{d}{dt}\Phi(z(t))=-\|\dot z(t)\|^2\quad\text{q.t.p.}
\end{equation}
A teoria não suave de Kurdyka--\L{}ojasiewicz para funções definíveis, combinada com a boa colocação de sistemas dinâmicos projetados com campo Lipschitz em conjuntos convexos fechados, fornece comprimento finito e convergência a um único ponto crítico projetado \cite{BDL,BDLS,Brogliato,Nagurney}.

\section{Minimalidade exata em hiperárvores lineares}
Seja $m_{*,HT}^{\rm mult}$ o menor número de cláusulas em um componente da classe HT para o qual existe um conjunto de condições iniciais de medida de Lebesgue positiva cujo limite arredondado viola ao menos uma cláusula.

\begin{theorem}[Obstrução mínima]
Para componentes conexos, lineares, 3-uniformes e Berge-acíclicos,
\[
m_{*,HT}^{\rm mult}=6.
\]
\end{theorem}

Fixe um equilíbrio $z^*$ de energia positiva e chame uma cláusula de positiva quando $P_c(z^*)>0$.

\begin{lemma}[Cláusula nula e coordenada livre]
Se $0<z_v^*<1$ e $P_c(z^*)=0$, então $\partial_vP_c(z^*)=0$.
\end{lemma}
\noindent A razão é que o fator de $v$ é estritamente positivo; como o produto é zero, algum outro fator é nulo e permanece após a derivação.

Considere uma componente conexa do sub-hipergrafo positivo com $t$ cláusulas. Berge-aciclicidade e 3-uniformidade dão $|V|=2t+1$. Se $L$ é o número de variáveis de grau um nessa componente,
\[
\sum_{v\in V}(\deg_+(v)-1)=3t-(2t+1)=t-1,
\]
portanto
\begin{equation}
L\ge t+2.
\end{equation}
Cada folha positiva deve estar no bordo superior, numa orientação em que seu fator positivo é $u_v$, e requer uma cláusula nula restauradora. Uma cláusula nula não pode restaurar duas folhas simultaneamente, pois dois fatores nulos anulam ambas as derivadas correspondentes. Se o componente completo tem $m$ cláusulas,
\[
m-t\ge L\ge t+2,
\]
e assim $m\ge2t+2$. Para $m\le5$, necessariamente $t=1$.

\subsection{Ramo unitário e folha monótona}
Se a única cláusula positiva é $C_0=y_1y_2y_3$, então $y_i^*=1$ e cada $y_i$ precisa de restauradora distinta. Removendo $C_0$ da árvore de incidência surgem três ramos não vazios. Para $m=4$, os tamanhos são $(1,1,1)$; para $m=5$, $(1,1,2)$. Logo sempre existe um ramo unitário. Após orientação,
\begin{equation}
A_y=(1-y)ab.
\end{equation}
As variáveis $a,b$ são folhas globais. No equilíbrio $y^*=1$,
\[
\partial_y\Phi(z^*)=1-a^*b^*\le0,
\]
portanto $a^*=b^*=1$. Além disso, globalmente,
\begin{equation}
\partial_a\Phi=(1-y)b\ge0,
\end{equation}
logo $\dot a\le0$. Qualquer trajetória que converja a esse equilíbrio deve ter $a(0)=1$. Sua bacia está contida em um hiperplano de bordo e tem medida ambiente zero. Como há apenas finitas escolhas possíveis de cláusula central e ramo unitário, a união relevante continua nula. Se a energia limite é zero, cada cláusula possui ao menos um fator de violação exatamente zero no endpoint Booleano satisfatório, então o arredondamento compatível satisfaz toda a fórmula. Assim, para $m\le5$, toda bacia ruim tem medida zero.

\section{A obstrução de seis cláusulas M6}
Considere o hiperarbusto assinado com vértices $0,\ldots,12$, arestas
\[
(0,1,2),(0,3,4),(1,5,6),(2,7,8),(3,9,10),(4,11,12),
\]
e triplas de sinais $(+++),(-++),(-++),(-++),(-++),(-++)$. Em coordenadas-fator,
\begin{align}
\Phi={}&yu_1u_2+(1-y)v_1v_2+(1-u_1)a_1b_1+(1-u_2)a_2b_2\
&+(1-v_1)c_1d_1+(1-v_2)c_2d_2.
\end{align}

\begin{proposition}[Família de equilíbrios positivos]
O conjunto
\[
\Mmot=\{u_1=u_2=v_1=v_2=1,\ 0<y<1,\ a_jb_j>y,\ c_jd_j>1-y\}
\]
é uma família de dimensão nove de equilíbrios projetados com $\Phi=1$ e consiste localmente em mínimos relativos não estritos.
\end{proposition}

A caixa aberta
\begin{equation}
U=\left\{\frac{31}{64}<y<\frac{33}{64},\quad \frac{15}{16}<u_j,v_j,a_j,b_j,c_j,d_j<1\right\}
\end{equation}
é capturada por $\Mmot$ em tempo finito. No envelope $7/16<y<9/16$, com os quatro fatores centrais pelo menos $15/16$ e as folhas acima de $7/8$, cada fator central não saturado tem derivada em coordenada-fator ao menos $13/64$, portanto velocidade no tempo original ao menos $13/256$. Assim ele alcança $1$ em tempo no máximo
\[
T\le \frac{1/16}{13/256}=\frac{16}{13}.
\]
Ao mesmo tempo,
\[
|\partial_y\Phi|\le\frac{31}{256},\qquad |\partial_a\Phi|\le\frac1{16},
\]
logo
\[
|\dot y|\le\frac{31}{1024},\qquad |\dot a|\le\frac1{64}.
\]
Os deslocamentos totais obedecem
\[
|\Delta y|\le\frac{31}{832}<\frac3{64},\qquad |\Delta a|\le\frac1{52}<\frac1{16},
\]
fechando o bootstrap. Uma vez saturados os quatro fatores centrais, todas as velocidades projetadas se anulam. Sob inicialização uniforme no cubo original, as coordenadas-fator são independentes uniformes em $[0,1]$, e
\begin{equation}
\mathbb P(U)=\frac1{32}\left(\frac1{16}\right)^{12}=2^{-53}=:p_0.
\end{equation}
Após arredondamento, exatamente uma das duas cláusulas centrais é violada, independentemente da convenção de empate em $y=1/2$.

\section{Ocorrência em fórmulas aleatórias esparsas}
Seja $M=\lfloor\alpha N\rfloor$ e $Q_N=8\binom N3$. Amostramos uniformemente entre os subconjuntos de $M$ cláusulas assinadas distintas. Seja $X_{M6}$ o número de componentes isolados na órbita de gauge de M6.

O motivo não assinado possui $|\mathrm{Aut}(M6)|=128$. A ação em embeddings injetivos é livre, e os $13$ flips de gauge produzem $2^{13}$ padrões assinados distintos por cópia estrutural. Para um suporte fixo de 13 variáveis,
\[
Q_{0,N}=8\binom{N-13}{3},
\]
e
\begin{equation}
\mathbb E X_{M6}=\frac{(N)_{13}}{128}\,2^{13}\frac{\binom{Q_{0,N}}{M-6}}{\binom{Q_N}{M}}.
\end{equation}
Como
\[
\frac{Q_{0,N}}{Q_N}=1-\frac{39}{N}+O(N^{-2}),
\]
segue que
\begin{equation}
\frac{\mathbb E X_{M6}}{N}\longrightarrow d(\alpha)=\frac{729}{64}\alpha^6e^{-39\alpha}.
\end{equation}
O modelo de cláusulas independentes possui a mesma constante líder.

\section{Concentração e cotas residuais}
Uma única troca de cláusula altera $X_{M6}$ em módulo no máximo quatro. No martingale de exposição de Doob do modelo de tamanho fixo, um acoplamento por transposição coloca as esperanças condicionais em um intervalo de comprimento no máximo quatro. Para $N$ suficientemente grande,
\begin{equation}
\mathbb P_F\!\left[X_{M6}<\frac34d(\alpha)N\right]\le \exp\!\left[-\frac{d(\alpha)^2}{512\alpha}N\right].
\end{equation}
Condicionado à fórmula, componentes isolados têm suportes disjuntos, portanto os eventos de captura dependem de blocos independentes da inicialização produto:
\[
Z_N\mid F\sim\mathrm{Binomial}(X_{M6},p_0).
\]
Chernoff fornece, no evento acima,
\[
\mathbb P_{x_0}\!\left[Z_N<\frac34p_0X_{M6}\mid F\right]\le \exp\!\left[-\frac{3p_0d(\alpha)}{128}N\right].
\]
Cada captura produz ao menos uma cláusula violada distinta, logo $\rho_{N,\mathrm{mult}}\ge Z_N/M$.

\begin{theorem}[Cotas residuais]
Para todo $\alpha>0$ fixo,
\[
\liminf_{N\to\infty}\mathbb E\rho_{N,\mathrm{mult}}\ge \frac{729}{2^{59}}\alpha^5e^{-39\alpha}>0.
\]
Além disso,
\[
\mathbb P\!\left(\rho_{N,\mathrm{mult}}\ge \frac{729}{2^{60}}\alpha^5e^{-39\alpha}\right)\to1.
\]
\end{theorem}

\section{Discussão e escopo}
O teorema determinístico é restrito a componentes conexos, lineares, 3-uniformes e Berge-acíclicos. Não afirmamos que M6 seja a menor obstrução entre componentes cíclicos ou não lineares. O resultado aleatório conta cópias isoladas do motivo certificado; não caracteriza todos os atratores de energia positiva.

A prova também separa duas questões frequentemente confundidas. O argumento KL estabelece convergência pontual de cada trajetória multilinear projetada a um único equilíbrio. A construção M6 mostra que convergência pontual não implica energia residual zero. A obstrução, portanto, não é a não convergência, mas a convergência para uma família de bordo não estrita com energia positiva.

\section{Conclusão}
Seis cláusulas são necessárias e suficientes, dentro de hiperárvores lineares 3-uniformes, para uma bacia ruim de medida positiva sob a dinâmica multilinear projetada de 3-SAT. O motivo mínimo ocorre com densidade linear em fórmulas aleatórias esparsas e produz cotas inferiores explícitas para o residual em esperança e com alta probabilidade assintótica. A combinação de geometria exata de pequenos componentes com contagem em estruturas aleatórias fornece um mecanismo concreto pelo qual uma extensão contínua exata nos vértices Booleanos pode conservar erro discreto não nulo após convergência.

\begin{thebibliography}{9}
\bibitem{BDL} J. Bolte, A. Daniilidis e A. Lewis, ``The \L{}ojasiewicz inequality for nonsmooth subanalytic functions with applications to subgradient dynamical systems'', \emph{SIAM J. Optim.} 17(4):1205--1223, 2007.
\bibitem{BDLS} J. Bolte, A. Daniilidis, A. Lewis e M. Shiota, ``Clarke subgradients of stratifiable functions'', \emph{SIAM J. Optim.} 18(2):556--572, 2007.
\bibitem{Brogliato} B. Brogliato, A. Daniilidis, C. Lemaréchal e J. Malick, ``Equivalence and complementarity issues in projected dynamical systems and differential inclusions'', \emph{Math. Programming} 107(3):317--335, 2006.
\bibitem{Nagurney} A. Nagurney e D. Zhang, \emph{Projected Dynamical Systems and Variational Inequalities with Applications}. Springer, 1996.
\end{thebibliography}
\end{document}